\documentclass[../../../include/open-logic-section]{subfiles}

\begin{document}

\olfileid{sth}{choice}{justifications}
\olsection{ఎంపికకు అంతర్గత సమర్థనపై పరిశీలనలు}

ఇక విస్తృతమైన ప్రశ్న: సుక్రమపరచడాన్నో, ఎంపికనో, లేదా అన్ని సమితుల పరిమాణాలను పోల్చవచ్చనే వాదననో---ఏది తీసుకున్నా తేడా లేదు---సమర్థించవచ్చా?

ఇక్కడ ఒక \emph{అంతర్గత} సమర్థన ప్రయత్నం ఉంది. \olref[z][story]{sec}లో స్థాయిక్రమం గురించి కొన్ని సూత్రాలను ప్రవేశపెట్టాం. వాటిలో ఒకదాన్ని మళ్లీ చెప్పాలి:
\begin{enumerate}
	\item[] \stagesacc. ఏ దశ $S$కైనా, ఆ దశ $S$కు \emph{ముందు} ఏర్పడిన సమితులను ఏవైనా తీసుకుంటే: ఆ సమితులే ఖచ్చితంగా మూలకాలుగా గల ఒక సమితి $S$ దశలో ఏర్పడుతుంది. ఆ దశ~$S$లో మరేదీ ఏర్పడదు.
\end{enumerate}
నిజానికి ఇలాంటి సూత్రం ద్వారా ఎంపిక స్వయంసిద్ధాన్ని సమర్థించవచ్చని పలువురు రచయితలు సూచించారు. ఆ పద్ధతిని సంక్షిప్తంగా వివరిద్దాం.

మొదట ఎంపికకు \emph{మరో} సమానార్థకాన్ని ఇచ్చే చిన్న ఫలితంతో మొదలెడదాం:

\begin{thm}[$\ZF$లో]\ollabel{choiceset}
$A$లోని \tetoken{మూలకాలు}{!!{element}s} జతలవారీ విచ్ఛిన్నాలై, ఖాళీ కానివైతే, ప్రతి $x \in A$కు $C
\cap x$ ఏకమూలక సమితి అయ్యే $C$ ఉంది అనే సూత్రానికి ఎంపిక సమానం. (అలాంటి $C$ను $A$ కోసం \emph{ఎంపిక సమితి} అంటాం.)
\end{thm}

ఈ ఫలితం నిరూపణ నేరుగా వస్తుంది; దాన్ని పాఠకుల అభ్యాసానికి వదిలేస్తాం.

\begin{prob}
\olref[sth][choice][justifications]{choiceset}ను నిరూపించండి. కష్టమైతే \cite[pp.~242--3]{Potter2004}లో నిరూపణ చూడవచ్చు.
\end{prob}

ముఖ్యమైన విషయం ఏమిటంటే $A$ కోసం ఎంపిక సమితి, $A$ కోసం ఎంపిక ప్రమేయపు వ్యాప్తి సమితియే. కాబట్టి ఎంపికను సమర్థించాలంటే ఎంపిక సమితుల అస్తిత్వం గురించిన దాని సమానార్థక రూపాన్ని సమర్థించడానికి ప్రయత్నించవచ్చు. ఇప్పుడు అదే చేద్దాం.

$A$లోని \tetoken{మూలకాలు}{!!{element}s} జతలవారీ విచ్ఛిన్నాలై, ఖాళీ కానివని అనుకుందాం. \stageshier{} ప్రకారం (\olref[z][story]{sec} చూడండి) $A$ ఏదో దశ~$S$లో ఏర్పడుతుంది. $\bigcup A$లోని \tetoken{మూలకాలు}{!!{element}s} అన్నీ $S$ దశకు ముందే అందుబాటులో ఉంటాయని గమనించండి. ఇప్పుడు \stagesacc{} ప్రకారం~$S$కు ముందు ఏర్పడిన \emph{ఏవైనా} సమితులకు, వాటినే ఖచ్చితంగా మూలకాలుగా గల సమితి ఏర్పడుతుంది. మరొక విధంగా చెప్పాలంటే: ముందే అందుబాటులో ఉన్న సమితుల ప్రతి \emph{సాధ్యమైన} సేకరణ~$S$లో ఉంటుంది. కానీ~$A$ కోసం ఎంపిక సమితి అయ్యే వస్తువులను ఎంచుకోవడం ఖచ్చితంగా \emph{సాధ్యం}; అది $\bigcup A$ యొక్క ఒక నిర్దిష్ట ఉపసమితి మాత్రమే. కనుక అవసరమైన ఎంపిక సమితి ఉంది.

ఇది అంతర్గత కారణాలతో ఎంపికను సమర్థించే \emph{చాలా} సంక్షిప్త ప్రయత్నం. ఈ ఆలోచనను ఇంకా పరిశీలించాలనుకుంటే, దాని గురించి పాటర్ ఇచ్చిన చక్కని వివరణను (\citeyear[\S14.8]{Potter2004}) చదవండి.

\end{document}
